\begin{problem}[Warmup]{(a)}
    What properties does $\mathbb{Z}[x]$ lack that prevent us from proving that it is a UFD just as we did for $\mathbb{Q}[x]$?
\end{problem}
\unsolved

\begin{problem}[Warmup]{(b)}
    Determine the content of \[30x^4 - 12x^2 + 42x - 54 \text{ and } 49x^3 + 70x^2 - 14.\] What is the content of the product of these two polynomials?
\end{problem}
\unsolved

\begin{problem}[Warmup]{(c)}
    Why would it be silly to try to talk about the content of polynomials from $\mathbb{Q}[x]$?
\end{problem}
\unsolved

\begin{problem}[Warmup]{(d)}
    Factor the following polynomials completely into irreducibles in $\mathbb{Z}[x]$; do they have the same irreducibles as factors in $\mathbb{Q}[x]$?
    \begin{enumerate}
        \item[(a)] $6x^3 - 6$.
        \item[(b)] $3x^4 - 6x$.
        \item[(c)] $5x^4 - x^3 - 15x^2 - 7x + 2$.
    \end{enumerate}
\end{problem}
\unsolved

\begin{problem}[Warmup]{(e)}
    Give examples of the following (or say why they don't exist):
    \begin{enumerate}
        \item[(a)] A PID that isn't a UFD.
        \item[(b)] A UFD that isn't a PID.
    \end{enumerate}
\end{problem}
\unsolved

\begin{problem}[Warmup]{(f)}
    Give examples of the following polynomials from $\mathbb{Z}[x]$ (or say why they don't exist):
    \begin{enumerate}
        \item[(a)] An irreducible polynomial of degree 4.
        \item[(b)] An irreducible polynomial that isn't prime.
        \item[(c)] An irreducible polynomial of degree 0.
        \item[(d)] A unit.
        \item[(e)] A polynomial that can't be non-trivially factored in $\mathbb{Z}[x]$, but can be non-trivially factored in $\mathbb{Q}[x]$.
        \item[(f)] A polynomial that can't be non-trivially factored in $\mathbb{Z}[x]$, but can be non-trivially factored in $\mathbb{R}[x]$.
        \item[(g)] A polynomial that can't be non-trivially factored in $\mathbb{Q}[x]$, but can be non-trivially factored in $\mathbb{Z}[x]$.
    \end{enumerate}
\end{problem}
\unsolved

\begin{problem}{1}
    Suppose that $f \in \mathbb{Z}[x]$ and the coefficient on its highest power of $x$ is
    1 (we say that such a polynomial is called \textbf{monic}).
    \begin{enumerate}
        \item[(a)] Why is a monic polynomial necessarily primitive?
        \item[(b)] Give an example of a primitive polynomial in $\mathbb{Z}[x]$ that isn't
        monic.
    \end{enumerate}
\end{problem}
\unsolved

\begin{problem}{2}
    Suppose that $f \in \mathbb{Z}[x]$ is monic (see Exercise 1).
    \begin{enumerate}
        \item[(a)] Prove that all rational roots of $f$ are integers.
        \item[(b)] Prove that all integer roots of $f$ divide its constant term.
        \item[(c)] Give examples of primitive polynomials for which parts a and b fail.
    \end{enumerate}
\end{problem}
\unsolved

\begin{problem}{3}
    Use Exercise 2 and the Root Theorem 4.3 to show that $x^3 + 2x + 7$ is prime in
    $\mathbb{Z}[x]$.
\end{problem}
\unsolved

\begin{problem}{4}
    Consider the ring $\mathbb{Z}[i][x]$ of polynomials with coefficients from the
    Gaussian integers $\mathbb{Z}[i]$. Argue that this ring is not a PID, by an argument
    analogous to that for $\mathbb{Z}$. (The argument in the text that $\mathbb{Z}[x]$ is
    a UFD can actually be generalized to this case too, but we will not do this here.)
\end{problem}
\unsolved

\begin{problem}{5}
    Prove that in $\mathbb{Z}[x]$, $\langle 3x + 1, x + 1 \rangle = \langle 2, x - 1 \rangle$.
    Show that this ideal is not principal.
\end{problem}
\unsolved
